\section*{Bab \ref{ch:g9:earth-resources} --- Sumber Daya dari Bumi}

\begin{solution}{exo:g9:earth-resources:1}
(a) \omterm{def:g9:earth-resources:aggregate}{Agregat}. (b) \omterm{def:g9:earth-resources:reserve}{Cadangan}. (c) \omterm{def:g9:earth-resources:ore-deposit}{Endapan bijih}. (d) \omterm{def:g9:earth-resources:ore-deposit}{Faktor konsentrasi}.
\end{solution}

\begin{solution}{exo:g9:earth-resources:2}
(a) Di dalam \omterm{def:g3:stones-pebbles-sand:pebble}{kerikil} \omterm{def:g1:landscapes-around-us:water}{sungai}. (b) Di dalam retakan, dari air panas. (c) Di
dalam \omterm{def:g7:volcanoes:magma}{magma}. (d) Inilah langkah pertama pembentukan endapan \omterm{def:g3:stones-pebbles-sand:pebble}{kerikil}
\omterm{def:g1:landscapes-around-us:water}{sungai}: emas yang terbebaskan kemudian dibawa oleh \omterm{def:g1:landscapes-around-us:water}{sungai}.
\end{solution}

\begin{solution}{exo:g9:earth-resources:3}
Urutannya: (b), (d), (a), (e), (c).
\end{solution}

\begin{solution}{exo:g9:earth-resources:4}
(a) Tidak terbarukan: batu bara memerlukan jutaan tahun untuk terbentuk.
(b) Terbarukan: \omterm{def:g4:clouds-rain-snow:precipitation}{hujan} mengisi kembali apa yang diambil. (c) Tidak
terbarukan: \omterm{def:g9:earth-resources:ore-deposit}{endapan bijih} terbentuk selama waktu yang sangat panjang.
(d) Tidak terbarukan dengan laju itu: \omterm{def:g6:water-underground:water-table}{muka air tanah} turun dan
\omterm{def:g6:water-underground:spring}{sumur-sumur} akhirnya kering.
\end{solution}

\begin{solution}{exo:g9:earth-resources:5}
(a) $4 \div 0.007 \approx 571$, sekitar $10^3$ (di antara $10^2$ dan
$10^3$, lebih dekat ke $10^3$). (b) $0.5 \div 0.0002 = 2500$, sekitar
$10^3$. (c) $20
\div 5.6 \approx 3.6$, sekitar 4: di antara $10^0$ dan $10^1$.
(d) Besi sudah sangat melimpah di \omterm{def:g6:inside-the-earth:layers}{kerak} bumi: bijih yang hanya beberapa
kali lebih kaya daripada \omterm{def:g3:stones-pebbles-sand:rock}{batuan} biasa sudah cukup.
\end{solution}

\begin{solution}{exo:g9:earth-resources:6}
(a) $19.3 \div 2.5 \approx 7.7$ kali. (b) Emas: $19.3 \times 0.01 =
\qty{0.193}{g}$; \omterm{def:g3:stones-pebbles-sand:rock}{batuan}: $2.5 \times 0.01 = \qty{0.025}{g}$. (c) Ketika
seluruh dasar \omterm{def:g1:landscapes-around-us:water}{sungai} bergerak, butir-butir emas yang berat tenggelam di
antara \omterm{def:g3:stones-pebbles-sand:sand}{pasir} dan \omterm{def:g3:stones-pebbles-sand:pebble}{kerikil} yang lebih ringan, sampai ke \omterm{def:g3:stones-pebbles-sand:rock}{batuan dasar} yang
padat, tempat butir-butir itu berhenti.
\end{solution}

\begin{solution}{exo:g9:earth-resources:7}
(a) Emas, timah, dan kromium. (b) Di antara $10^1$ dan $10^2$, dekat
dengan $10^2$. (c) $3000 \div 3.7 \approx 800$ kali. (d) Kromium (atau
timah, atau emas): makin tinggi \omterm{def:g9:earth-resources:ore-deposit}{faktor konsentrasinya}, makin jarang
tempat alam telah mengonsentrasikan logam itu dengan cukup.
\end{solution}

\begin{solution}{exo:g9:earth-resources:8}
(a) $18 \times 365 = \qty{6570}{kg}$, sekitar \qty{6.6}{t}. (b) $5 \times
10^{10}$ ton. (c) $25000 \div 18 \approx 1389$ hari, hampir empat tahun.
(d) \omterm{def:g3:stones-pebbles-sand:sand}{Pasir} yang digali dari \omterm{def:g1:landscapes-around-us:water}{sungai} tidak lagi terbawa ke hilir sampai ke
pantai; ombak terus membawa \omterm{def:g3:stones-pebbles-sand:sand}{pasir} pergi dari pantai, tetapi \omterm{def:g3:stones-pebbles-sand:sand}{pasir} yang
datang untuk menggantikannya berkurang.
\end{solution}

\begin{solution}{exo:g9:earth-resources:9}
(a) Dari tumbuhan dan hewan renik di \omterm{def:g2:water-on-earth:ocean}{laut}, seperti diatom, yang terkubur
di dalam \omterm{def:g3:stones-pebbles-sand:sand}{lumpur} dasar \omterm{def:g2:water-on-earth:ocean}{laut}. (b) Antara sekitar 60 dan
\qty{120}{\degreeCelsius}. (c) Keduanya lebih ringan daripada air, yang
mengisi \omterm{def:g6:water-underground:porosity}{pori-pori} \omterm{def:g3:stones-pebbles-sand:rock}{batuan}. (d) Lapisan yang kedap, seperti serpih; keduanya
berkumpul di bawahnya, sering kali di puncak lipatan yang melengkung ke
atas, gas di atas minyak.
\end{solution}

\begin{solution}{exo:g9:earth-resources:10}
(a) $2.9 \times 10^7$ ton. (b) $2.9 \times 10^7 \div 100 = 2.9 \times
10^5$ ton setahun, kira-kira. (c) Bijihnya 100 kali lebih berat daripada
tembaga yang dikandungnya: $2.9 \times 10^7 \times 100 = 2.9 \times 10^9$
ton. (d) Setiap ton bijih hanya menghasilkan \qty{10}{kg} tembaga,
sehingga miliaran ton \omterm{def:g3:stones-pebbles-sand:rock}{batuan} harus digali.
\end{solution}

\begin{solution}{exo:g9:earth-resources:11}
(a) Dengan 23~juta ton setahun, \omterm{def:g9:earth-resources:reserve}{cadangan} itu akan bertahan sekitar
$280 \div 23
\approx 12$ tahun, sampai awal tahun 1980-an; namun, 697~juta ton telah
ditambang, jauh lebih banyak daripada \omterm{def:g9:earth-resources:reserve}{cadangan} itu.
(b) $980 \div 280 = 3.5$ kali. (c) $980 \div
23 \approx 43$ tahun. (d) Misalnya: endapan baru ditemukan; teknik yang
lebih baik atau harga yang lebih tinggi membuat endapan yang lebih miskin
layak ditambang.
\end{solution}

\begin{solution}{exo:g9:earth-resources:12}
(a) $5 \times 10^1$ tahun. (b) $2 \times 10^3 \div 50 = 40$ kali. (c) $3
\times 10^8 \div (5 \times 10^1) = 6 \times 10^6$ kali. (d) Batu bara
memerlukan waktu jutaan kali lebih lama untuk terbentuk daripada waktu
yang kita perlukan untuk memakainya: alam tidak memperbaruinya pada skala
waktu kita.
\end{solution}

\begin{solution}{pb:g9:earth-resources:1}
\textbf{1.} $100 \times 1000 = \num{100000}$~kg $= 1 \times 10^5$~kg.

\textbf{2.} $2000 \times 300 = \num{600000}$~t $= 6 \times 10^5$~t.

\textbf{3.} $6 \times 10^5 \div 10 = 6 \times 10^4$~t per tahun.

\textbf{4.} $0.40 \times 6 \times 10^4 = 2.4 \times 10^4$~t \omterm{def:g3:stones-pebbles-sand:sand}{pasir} per
tahun.

\textbf{5.} $1200 \times 500 \times 5 = 3 \times 10^6$~m$^3$.

\textbf{6.} Dua pertiganya: $2 \times 10^6$~m$^3$.

\textbf{7.} $2 \times 10^6 \times 1.6 = 3.2 \times 10^6$~t.

\textbf{8.} $\num{24000} + \num{136000} = \num{160000}$~t $= 1.6 \times
10^5$~t per tahun.

\textbf{9.} $\num{160000} \div 2000 = 80$ kali lebih banyak.

\textbf{10.} Tidak: lubang galian itu dikosongkan 80 kali lebih cepat
daripada \omterm{def:g1:landscapes-around-us:water}{sungai} memperbarui \omterm{def:g3:stones-pebbles-sand:sand}{pasirnya}. Pada skala kota itu, \omterm{def:g3:stones-pebbles-sand:sand}{pasir} ini
adalah sumber daya yang terbatas.

\textbf{11.} $\num{24000} \div 4 = \num{6000}$~t \omterm{def:g3:stones-pebbles-sand:sand}{pasir} setahun.

\textbf{12.} $3.2 \times 10^6 \div (1.6 \times 10^5) = 20$: \textbf{20 tahun}.
\end{solution}
